\section{Les tabliers des notes de Grothendieck}\label{sec:grdecks}

Trois tabliers reviennent dans les notes de travail de Grothendieck. Voici ce qu'en disent les trois arches du \S\ref{sec:morearches} (Lean: \lean{GrothendieckDecks.lean}).
\begin{itemize}
\item La catégorie simpliciale $\Delta$, la première catégorie test de \emph{Pursuing Stacks} (dossiers 108 et 118 du fonds): les ensembles simpliciaux forment un topos à deux valeurs et connexe, puisqu'il existe une application constante entre deux objets quelconques de $\Delta$, et ce n'est pas le topos de préfaisceaux d'un préordre, car les deux sommets de $[1]$ sont deux flèches $[0]\to[1]$ (\lean{simplexCategory_twoValued}, \lean{simplexCategory_isConnected}, \lean{simplexCategory_not_spatial}).
\item Un groupe $G$ vu comme groupoïde à un objet, dont le topos de préfaisceaux est $\B G$ (dossier 161-6, p.~22, où Grothendieck ramène une question à ce cas): à deux valeurs, comme le note \cite[IV~8.4.3]{sga4} pour tout monoïde, connexe, et topos de préfaisceaux d'un préordre si et seulement si $G$ est trivial (\lean{singleObj_twoValued}, \lean{singleObj_isConnected}, \lean{singleObj_spatial_iff}).
\item Un ensemble préordonné $I$, dont le topos de préfaisceaux $\widehat I$ est $\mathrm{Top}(I)$ (dossier 37, p.~48; dossier 161-3, p.~26): à deux valeurs si et seulement si deux éléments quelconques sont comparables dans les deux sens, c'est-à-dire si et seulement si son quotient ordonné a au plus un point (\lean{preorder_twoValued_iff}).
\end{itemize}
Ce sont des corollaires immédiats; nous les consignons parce que ce sont les premières entrées du registre tirées du fonds.
